% !PRINTED 740
\weberpage{740}

%<<CONT
$a_{1,1}$ is this odd number. Then we determine a
second solution:
\[
0, a_{2,2}, a_{3,2}, \ldots, a_{\nu-1,2},
\]
in which again an odd number, say $a_{2,2}$, occurs. Then we
determine a third solution:
\[
0, 0, a_{3,3}, \ldots, a_{\nu-1,3},
\]
and we continue setting to zero as long as the number of the
unknowns still remaining is greater than the number of the
equations.

In the last system, therefore, $\nu'-1$ zeros will occur, so
that $\nu'$ unknowns remain.

We may take these rows of numbers $a_{\nu,k}$ as the system (20),
for in it the determinant
\weq{21}{\Sigma \pm a_{1,1}\, a_{2,2} \ldots a_{\nu',\nu'}}
reduces to the diagonal term $a_{1,1}\, a_{2,2} \ldots a_{\nu',\nu'}$, whose factors are all
 odd.

If now the condition in statement 2 is satisfied, the matrix
(20) can be completed, by the addition of $\nu'-1$ rows, to a
determinant of $\nu-1$ rows which likewise has an odd integral
value:
\weq{22}{a = \Sigma \pm a_{1,1}\, a_{2,2} \ldots a_{\nu',\nu'}\, a_{\nu'+1,\nu'+1} \ldots a_{\nu-1,\nu-1}.}

One need only, when the determinant (21) is assumed to be odd,
set nothing but zeros for the elements in the added rows, except
those on the diagonal,
 $a_{\nu'+1,\nu'+1} \ldots a_{\nu-1,\nu-1}$, which may be taken equal to 1. Then $a$
receives the same value as the determinant
 (21).

But if we now form the product $a M$ according to the
multiplication theorem for determinants, there results a
determinant in which the $\nu-1$ sums
\[
\begin{aligned}
& \sum\limits_{1,\, \nu-1}^{i} a_{i,s}\, m_{k,i}, \qquad \sum\limits_{1,\, \nu-1}^{i} a_{i,s}\, M_{k,i} \\
& k = 1, 2, \ldots \nu'-1, \qquad k = 0, 1, 2 \ldots \nu'-1
\end{aligned}
\]
form a row for each value $s = 1, 2, \ldots, \nu'$, and in which
therefore $\nu'$ rows occur whose elements are all divisible by
 2.
